\chapter{Symbols, Terminology, and Glossary}
\label{sec:p14-symbols-terminology-and-glossary}
\index{glossary}\index{terminology}\index{notation!lookup}

A notation lookup asks what a written sign denotes in its present context.
A terminology lookup asks which definition a word invokes. Both require the
language, the type of its operands and the local convention: a turnstile can
express satisfaction or consequence, and an independence condition can concern
events, preferences or potential outcomes. This chapter supplies a selected
cross-domain vocabulary and links to its exact home definitions and results.
The symbol registry in \cref{sec:p01-symbol-registry}, the model-theory
subtable in \cref{sec:p03-model-symbols}, and the named-theorem index in
\cref{sec:p14-theorem-index} provide the complementary retrieval surfaces.
The source-specific grammars and hypotheses at those destinations govern;
the summaries here neither add axioms nor extend a theorem's scope
\cite{Enderton2001,HalpernUncertaintyChapter2Extract,Pearl2009CausalOverview,
Rubinstein2015LectureNotes}.

\section{Symbols by spelling, role, and source convention}
\label{sec:p14-glossary-symbols}

\begin{hbmethod}{def:p14-glossary-symbol-lookup}{Reading a symbol in context}
Identify the chapter's named presentation; read the declaration of the
expression's grammatical category and operand types; distinguish an
object-language constructor from a metalanguage relation or abbreviation;
then consult the scoped registry row and its first-use destination. When a
source uses a familiar glyph differently, retain that source's declared
syntax and read the typed local definition. A spelling match alone does not
identify a mathematical operation. Parentheses, subscripts and the operand
types are part of the lookup, not optional decoration.
\end{hbmethod}

The following rows give high-use spellings. A macro spelling is an authoring
lookup; the printed expression and the local declaration determine its
mathematical meaning. The generated registry contains the individual scope,
precedence and first-use records, rather than assigning a meaning to every
letter of a font family. The historical notation table in
\cref{tab:p01-notation-variants} records named editions, not an invariant
notation shared by all authors of a tradition.

\begingroup\small
\begin{longtable}{@{}p{.19\textwidth}p{.23\textwidth}p{.48\textwidth}@{}}
\caption{Common spellings, typed roles and defining destinations.}
\label{tab:p14-glossary-spellings}\\
\toprule
Printed form & TeX spelling & Role and source convention \\
\midrule
\endfirsthead
\caption[]{Common spellings, typed roles and defining destinations (continued)}\\
\toprule
Printed form & TeX spelling & Role and source convention \\
\midrule
\endhead
\midrule\multicolumn{3}{r}{\emph{continued on next page}}\\
\endfoot
\endlastfoot
\(\neg,\land,\lor,\to\) &
\texttt{\detokenize{\neg \land \lor \to}} &
Object-language connectives when admitted by the active grammar; an
intuitionistic presentation defines negation by implication into falsum.
\cref{sec:prop-syntax-semantics-01,def:p08-ipl-completeness-setting}. \\
\(\leftrightarrow\) & \texttt{\detokenize{\leftrightarrow}} &
Object-language biconditional, primitive or expanded as declared; not the
assertion that two expressions are definitionally identical.
\cref{sec:p01-binding-precedence,sec:prop-syntax-semantics-01}. \\
\(\Leftrightarrow\) & \texttt{\detokenize{\Leftrightarrow}} &
Editorial metalinguistic formula equivalence. The source CPL presentation
instead explicitly retains \(\equiv\) for truth-value equivalence.
\cref{def:p02-uniform-substitution,thm:p02-replacement}. \\
\(\defeq\) & \texttt{\detokenize{\defeq}} &
Definitional abbreviation; it declares a meaning, not a theorem establishing
existence or equality independently of that declaration.
\cref{sec:p01-symbol-registry,sec:levels-02}. \\
\(\Gamma\derives_{\mathcal C}\varphi\) &
\texttt{\detokenize{\derives}} &
Derivability from open assumptions in the specified calculus.
\cref{sec:levels-03}; the rule sheets supply the exact licensed rules. \\
\(\mathcal M,s\entails\varphi\),
\(\Gamma\entails\varphi\) & \texttt{\detokenize{\entails}} &
Respectively satisfaction at supplied semantic data and consequence over
all permitted such data. The left operand distinguishes these relations.
\cref{sec:fol-semantics-02,sec:levels-03}. \\
\(\Gamma\Rightarrow\Delta\) & \texttt{\detokenize{\Rightarrow}} &
Sequent punctuation with the active sheet's context convention. Numeric
operands in the t-norm chapter instead give its residuum operation.
\cref{sheet:lk,sheet:lj,def:p08-tnorm-residuum}. \\
\(\forall x,\exists x\) &
\texttt{\detokenize{\forall \exists}} &
First-order binders. The categorical forms \(\forall(l,m)\) and
\(\exists(l,m)\) instead combine two literals into a complete sentence;
they bind no written individual variable.
\cref{sec:fol-syntax-semantics-02,def:p02-syllogistic-language}. \\
\(\FV(\varphi)\), \(\varphi[t/x]\) &
\texttt{\detokenize{\FV}}, \texttt{[t/x]} &
Free-variable set and capture-avoiding substitution for free occurrences.
The operand category fixes formula, term or context action.
\cref{sec:fol-syntax-semantics-03,def:p02-fol-capture-avoiding-substitution}. \\
\(\operatorname{Mod}(T)\), \(\operatorname{Th}(\mathcal M)\) &
\texttt{\detokenize{\Mod \Th}} &
Structures satisfying the theory; sentences true in the structure. Fix the
signature, including any parameter expansion.
\cref{sec:levels-04,def:parameter-expansion}. \\
\(\mathcal M\equiv\mathcal N\),
\(\mathcal M\preccurlyeq\mathcal N\),
\(\mathcal M\cong\mathcal N\) &
\texttt{\detokenize{\equiv \preccurlyeq \cong}} &
Elementary equivalence, elementary substructure and isomorphism; the
model-theory subtable keeps these distinct from source formula equivalence
and computability's Kleene equality.
\cref{tab:p03-model-symbols,lem:iso-invariance}. \\
\(\operatorname{Diag}(\mathcal M)\),
\(\operatorname{ElDiag}(\mathcal M)\) &
\texttt{\detokenize{\operatorname{Diag}}},
\texttt{\detokenize{\operatorname{ElDiag}}} &
Atomic and negated atomic parameter sentences, versus all parameter
sentences true in the expansion. \cref{tab:p03-model-symbols}. \\
\(S_n(T)\) & \texttt{S\_n(T)} &
Space of complete \(n\)-types consistent with \(T\), with the declared
clopen topology; not the tuples realized in one structure.
\cref{def:type-space,thm:type-space-compact}. \\
\(\mathbb P\) & \texttt{\detokenize{\mathbb P}} &
A probability measure when its operands are events, but a forcing poset in
the set-theoretic forcing presentation. Consult the typed registry rows;
neither role follows from blackboard-bold type alone.
\cref{def:p15-finite-probability-space,
def:p06-forcing-order-orientation,sec:p01-letter-conventions}. \\
\(w\forces A\) & \texttt{\detokenize{\forces}} &
Intuitionistic forcing at a world with persistent atomic information.
Set-theoretic forcing has a condition, names and a ground-model/poset
declaration instead. \cref{def:p08-ipl-completeness-semantics,
sec:constructibility-forcing-02}. \\
\(\Box,\Diamond\) & \texttt{\detokenize{\Box \Diamond}} &
Modal operators in a declared modal language; modal necessity is not a
numerical necessity measure. Source-local relevance abbreviations have
their own role. \cref{sec:p14-modal-correspondence,
def:p14-evidence-uncertainty-data,sec:p01-symbol-registry}. \\
\(\otimes,\multimap\) &
\texttt{\detokenize{\otimes \multimap}} &
Linear-logic tensor and linear implication under their sheet. Inquisitive
tensor instead splits an information state into a union of two supporting
states. \cref{def:p08-inq-tensor,def:p08-cll-language,
sec:p01-binding-precedence}. \\
\(\bot,\top\) & \texttt{\detokenize{\bottom \topv}} &
Falsum and verum in the ordinary registered role. Classical linear logic
assigns its declared nullary-unit roles locally, rather than importing
ordinary truth-functional readings. \cref{def:p08-cll-language}. \\
\(\mathbb P(A\mid B)\) & \texttt{\detokenize{\mid}} &
Conditioning on a positive-probability event. A bar followed by
\(\operatorname{do}(V_I=x)\) denotes the separately defined intervention
operation, not an event being conditioned on.
\cref{def:p15-finite-conditioning,def:p16-causal-intervention}. \\
\(\underline P,\overline P\) &
\texttt{\detokenize{\underline P \overline P}} &
Event envelopes over retained laws in the finite comparison, or lower and
upper previsions on gambles in the lower-prevision chapter; their domains
must be read. \cref{def:p14-evidence-uncertainty-data,def:p15-ip-prices}. \\
\(\succeq,\succ,\sim\) &
\texttt{\detokenize{\succeq \succ \sim}} &
Weak preference, strict preference and indifference in the finite preference
chapter. Indifference requires two affirmative weak comparisons, not a lack
of information. \cref{def:p17-weak-preference}. \\
\(s,\Gamma\) & \texttt{s}, \texttt{\detokenize{\Gamma}} &
First-order assignment and formula context under global conventions;
the best-first-search instance declares an initial graph state and goal
set instead. \cref{sec:fol-semantics-01,def:p12-hsearch-instance}. \\
\bottomrule
\end{longtable}
\endgroup

Letters such as \(\alpha,\beta\) likewise have no untyped universal role:
they may be ordinals, types, test levels or probabilities under different
declarations. In \cref{def:p15-statistical-experiment}, a test's level is an
upper bound on size; \cref{def:p15-confidence-procedure} explicitly notes
Neyman's different use of \(\alpha\) for a confidence coefficient.
Read the source's definition before transferring a familiar formula.

\section{Formula grammar, precedence, binding, and substitutions}
\label{sec:p14-glossary-grammar}

\subsection{Grammar before precedence}

A signature declares nonlogical symbols and arities; a grammar determines
terms and formulas. Precedence merely abbreviates an already admitted parse.
It cannot make a string belong to a language whose constructors do not include
the displayed symbol. The fully parenthesized first-order grammar is in
\cref{sec:fol-syntax-semantics-01,sec:fol-syntax-semantics-02}; the
propositional readability result is \cref{thm:prop-unique-readability}.
For example, the categorical fragment of
\cref{def:p02-syllogistic-language} admits neither arbitrary conjunctions
nor nested variable-binding quantifiers, although its printed categorical
operators are also called universal and existential.

The editorial ladder in \cref{tab:p01-precedence} has smaller ranks bind
more tightly: application first; atomic relations next; unary operators and
quantifiers; temporal until; conjunction; disjunction; implication;
biconditional. Its linear exponential is at the separately declared rank
between atomic relations and the other unary operators. Conjunction and
disjunction associate to the left, implication to the right, and repeated
biconditionals and until operators require parentheses. This is a reading
convention, not permission to identify different parse trees in every logic.
The source's explicitly local convention takes precedence within a named
source presentation \cite{WhiteheadRussell1910,Lukasiewicz1929,
HilbertAckermann1928,Gentzen1934a}.

In particular, the active source CPL grammar of
\cref{sec:prop-syntax-semantics-01} right-associates repeated conjunction,
disjunction and implication. Its \(P\land Q\land R\) abbreviates
\(P\land(Q\land R)\), whereas the editorial left-associative convention
abbreviates \((P\land Q)\land R\). The source also fixes its finite
indexed conjunctions and disjunctions as right-associated strings. Logical
equivalence of particular bracketings in a particular semantics does not
make their formula strings identical or determine another system's
associativity convention.

Thus the registered abbreviations read
\[
 \neg P\land Q\to R\quad\text{as}\quad
 ((\neg P)\land Q)\to R,
 \qquad
 P\to Q\to R\quad\text{as}\quad P\to(Q\to R).
\]
Under the same convention,
\(\forall x\,P(x)\to Q(x)\) means
\((\forall x\,P(x))\to Q(x)\), not
\(\forall x\,(P(x)\to Q(x))\). If the latter is intended, write its
parentheses. Untyped lambda application has a different explicit local
convention: \(MNP\) means \((MN)P\), and application binds more tightly
than abstraction; see \cref{def:p07-untyped-lambda-terms}.

\subsection{Binding belongs to occurrences}

An occurrence is free or bound relative to the parse tree. In
\(P(x)\land\exists x\,Q(x,y)\), the occurrence of \(x\) in the left
conjunct is free and the occurrence in the right conjunct is bound; \(y\)
is free. A sentence has no free variable occurrence. The recursive free-variable
and all-variable sets are different operations, as specified in
\cref{sec:fol-syntax-semantics-03}. For first-order satisfaction the
assignment cannot be suppressed merely because some other occurrences of
the same written variable are bound. Assignment independence is governed by
\cref{lem:fol-coincidence}, and the truth-in-a-structure convention for
sentences by \cref{sec:levels-04}.

Alpha-equivalence consistently changes bound names without changing the
binding pattern; it is not arbitrary replacement of a variable symbol.
First-order alpha-equivalence and its semantic invariance are at
\cref{def:p02-fol-alpha-equivalence,cor:alpha-invariance}. The lambda-term
relation in \cref{def:p07-alpha-substitution} has its own term grammar and
is not the same relation on the same set of objects.

\subsection{Four different replacement operations}

\begin{longtable}{@{}p{.23\textwidth}p{.42\textwidth}p{.25\textwidth}@{}}
\caption{Four different replacement operations}\\
\toprule
Operation & What is replaced; what is preserved & Exact destination \\
\midrule
\endfirsthead
\caption[]{Four different replacement operations (continued)}\\
\toprule
Operation & What is replaced; what is preserved & Exact destination \\
\midrule
\endhead
\midrule\multicolumn{3}{r}{\emph{continued on next page}}\\
\endfoot
\endlastfoot
Uniform propositional substitution & Every original occurrence of each
propositional variable receives its assigned formula simultaneously;
inserted formulas are not substituted again. &
\cref{def:p02-uniform-substitution,thm:p02-uniform-substitution}. \\
First-order term substitution & A term replaces free occurrences of one
individual variable; a conflicting binder is renamed before insertion. &
\cref{def:p02-fol-capture-avoiding-substitution,thm:fol-substitution}. \\
Selected subformula replacement & Only the specified original parse-tree
positions are replaced. Equivalence preservation requires the replacement
theorem's logical-equivalence hypothesis, not mere equisatisfiability. &
\cref{def:p02-selected-replacement,thm:p02-replacement,
def:p02-substitution-replacement-limits}. \\
Lambda substitution and beta contraction & Capture-avoiding term substitution
is used when a lambda redex contracts. Alpha-renaming, directed reduction and
symmetric conversion remain distinct. &
\cref{def:p07-alpha-substitution,def:p07-beta-conversion}. \\
\bottomrule
\end{longtable}

For distinct variables \(x,y\) and a binary relation symbol \(R\), the
first-order substitution \((\forall y\,R(x,y))[y/x]\) requires a fresh
binder replacing \(y\) before the term \(y\) is inserted. With the
fresh variable chosen by the definition denoted here by \(z\), the result
is \(\forall z\,R(y,z)\), not \(\forall y\,R(y,y)\).
The former leaves the inserted \(y\) free; the latter captures it.
The exact construction fixes an ordered supply of variables, so its literal
fresh name is determined there. An arbitrary fresh representative describes
the same alpha-class, not necessarily the same literal string.
The alpha-class/literal distinction is also maintained in the exact
cut-elimination construction at \cref{thm:p07-lj-exact-cut-elimination}.

For simultaneous propositional substitution sending \(P\) to \(Q\)
and \(Q\) to \(\neg Q\), \((P\land P)\) becomes \((Q\land Q)\),
not \((\neg Q\land\neg Q)\). Replacing only its left occurrence of
\(P\) by \(Q\) instead gives \((Q\land P)\).
These are the different operations of
\cref{def:p02-substitution-replacement-example}, not competing spellings
for one procedure.

\section{Terms with different meanings in different systems}
\label{sec:p14-glossary-overloads}

The adjective and the quantified objects usually resolve an overloaded term.
The following comparisons state both readings instead of forcing them into
one definition. A result proved for one row's meaning cannot be applied to
another merely because the same word occurs.

\begin{longtable}{@{}p{.15\textwidth}p{.37\textwidth}p{.38\textwidth}@{}}
\caption{Meaning changes that affect reasoning.}
\label{tab:p14-glossary-overloads}\\
\toprule
Term & One meaning & Distinct meaning and boundary \\
\midrule
\endfirsthead
\caption[]{Meaning changes that affect reasoning (continued)}\\
\toprule
Term & One meaning & Distinct meaning and boundary \\
\midrule
\endhead
\midrule\multicolumn{3}{r}{\emph{continued on next page}}\\
\endfoot
\endlastfoot
Admissible & A set of arguments is conflict-free and defends every member.
\cref{def:p18-defense-operator}. & An admissible search heuristic does not
exceed the optimum remaining cost. \cref{def:p12-hsearch-heuristics}.
Neither is rule admissibility in \cref{sec:p14-glossary}. \\
Argument & Premises and conclusion as an ordered pair.
\cref{sec:levels-05}. & An argument identifier in an abstract attack graph;
its internal premises and inference are not supplied by the graph.
\cref{def:p18-finite-framework}. \\
Belief & A modal operator's truth at accessible worlds, with factivity
depending on the system. \cref{entry:p08-kd45n}. & A mass-based numerical
belief function aggregates subset masses; a Bayesian posterior is another
specified probability law. \cref{def:p14-evidence-uncertainty-data,
def:p15-bayes-experiment}. \\
Coherence & Lower-prevision portfolio constraints include avoiding sure
loss. \cref{def:p15-ip-coherence}. & Psychological agreement among retrieved
considerations is a proposed confidence cue, not a proof of those portfolio
constraints or logical consistency. \cref{def:p19-r86-cue-standing,
sec:p19-r86-cues}. \\
Complete & A calculus derives every semantic consequence in the declared
scope. \cref{sec:levels-03}. & A theory decides every sentence; a model-complete
theory has only elementary embeddings between its models; a complete
argumentation extension contains everything it defends.
\cref{def:theory-completeness,def:p18-extension-semantics}. \\
Confidence & A judgment about a specified answer or event, interpreted on
its recorded scale. \cref{def:p19-calibration-law,
def:p19-r86-judgment-types}. & A confidence procedure has sampling coverage
for each fixed parameter; a credible set has conditional posterior content.
\cref{def:p15-confidence-procedure,def:p15-bayes-credible-set}. \\
Consistency & Nonderivability of falsum in a specified calculus; some
model-theory and abduction passages explicitly use satisfiability and state
the connecting metatheory. \cref{def:type-space,def:p16-abduction-models}. &
Heuristic consistency is an edgewise cost inequality; potential-outcome
consistency equates the observed response with the response under the
received treatment. \cref{def:p12-hsearch-heuristics,
def:p16-binary-potential-outcomes}. \\
Independence & Event or random-variable product factorization under a named
probability law. \cref{def:p15-finite-independence}. & Preference mixture
independence compares lotteries with a common remaining component.
\cref{def:p17-lottery-independence}. A causal exchangeability assumption is
probabilistic independence in a joint law including potential outcomes,
not a fact read directly from observed variables.
\cref{def:p16-binary-potential-outcomes}. \\
Model & A structure satisfying the specified sentences; bare syntax does
not make it a model of those sentences. \cref{sec:levels-04}. & A statistical
experiment specifies sampling laws; a structural causal model additionally
specifies equations and intervention meanings.
\cref{def:p15-statistical-experiment,def:p16-finite-scm}. \\
Normal & A beta-normal term has no beta redex. \cref{def:p07-beta-conversion}.
A formula normal form has its stated grammatical restriction. & A normal
random-effects model makes distributional assumptions, not a syntactic
normal-form assertion. \cref{def:p15-synthesis-random-model}.
Normalization of a probability mass means total mass one.
\cref{def:p15-finite-probability-space}. \\
Possibility & Modal possibility depends on accessibility and the selected
modal truth clauses. \cref{sec:p14-modal-correspondence}. & A numerical
possibility measure uses a maximum of singleton scores; it is generally
not an additive probability. \cref{def:p14-evidence-uncertainty-data,
thm:p14-evidence-measure-separation}. \\
Preference & A person's or decision model's weak comparison of alternatives.
\cref{def:p17-weak-preference}. & A preferential logical model's state ordering
selects minimal premise states for defeasible consequence; it does not
automatically represent utility or a measured person's choices.
\cref{def:p08-preferential-model}. \\
Resolution & A refutation method has the resolution theorem's syntactic
clause and substitution conditions. \cref{thm:fol-resolution}. & Forecast
resolution measures variation in outcome rates between judgment classes;
it is neither proof resolution nor the variance of confidence reports.
\cref{def:p19-calibration-law,thm:p19-calibration-decomposition}. \\
Support & A discrete law's support contains positive-mass outcomes.
\cref{def:p15-finite-probability-space}. & Inquisitive support is a relation
between a formula and an information state, recursively defined over
substates; it assigns no probability mass. \cref{def:p08-inq-support}.
Empirical support for a claim requires the evidence record.
\cref{def:p19-evidence-record}. \\
Type & A model-theoretic type is a set of formulas consistent with the
specified theory or parameter structure. \cref{def:type-space}. & A type
in a term calculus belongs to its type grammar and typing judgment; neither
is a type-I or type-II statistical error.
\cref{sec:p07-lambda-calculi-and-type-theories,
def:p15-statistical-experiment}. \\
Valid & A formula is true at every permitted semantic datum.
\cref{sec:levels-03}. & A valid \(p\)-value is superuniform under every null
law in its stated family; cue validity concerns a cue's empirical relation
to its target. \cref{def:p15-valid-p-value,def:p19-r86-cue-standing}. \\
\bottomrule
\end{longtable}

\section{Glossary of logic, inference, probability, causation, and decision}
\label{sec:p14-glossary-domains}
\subsection{Glossary}
\label{sec:p14-glossary}

This compact logical vocabulary is complemented by the cross-domain lookup
in \cref{sec:p14-symbols-terminology-and-glossary}, which distinguishes
source conventions and overloaded terms across reasoning disciplines.

\begin{description}[style=nextline,leftmargin=0pt]
\item[admissible rule] A rule whose conclusion is derivable after every
substitution making all premises derivable, whether or not the rule itself is
derivable.
\item[algebraic semantics] An interpretation of logical operations and
consequence in a declared class of algebras; see Chapter~\ref{sec:p10-algebraic-and-categorical-logic}.
\item[compactness] The property that semantic consequence from a set is already
consequence from a finite subset; first-order compactness is
Theorem~\ref{thm:fol-compactness}.
\item[complete theory] A consistent theory deciding every sentence of its
language; this differs from completeness of a calculus.
\item[congruence] An operation-compatible equivalence relation on an algebra;
see Definition~\ref{def:p10-algebra}.
\item[conservative extension] An extension proving no new sentences in the old
language beyond those already provable in the base theory.
\item[decidable set] A set with a total computable characteristic function.
\item[derivability] Existence of a derivation in a named calculus; never an
unsubscripted synonym for semantic consequence.
\item[effective procedure] A finitely specified method analyzed through a
declared formal model of computation.
\item[elementary embedding] An embedding preserving every first-order formula
with parameters.
\item[equisatisfiable] Having a model simultaneously in the sense that one
formula is satisfiable iff the other is; weaker than equivalence.
\item[filter] An upward-closed, finite-meet-closed subset of a lattice; see
Definition~\ref{def:p10-boolean-algebra}.
\item[finite validity] Truth in every finite structure of the fixed signature;
see Definition~\ref{def:p09-finite-validity}.
\item[frame] The relational component of a Kripke model, before a valuation is
added.
\item[free algebra] An algebra satisfying the extension universal property from
a chosen generating set.
\item[homomorphism] A map preserving every operation or relation specified by
the context.
\item[interpretation of theories] A translation of one theory into another
with a declared domain formula and preservation of axioms/proofs.
\item[invariant] A property unchanged under a specified class of transformations;
the transformation class is part of the assertion.
\item[Kripke model] A frame together with a valuation or domain/interpretation
data appropriate to the modal language.
\item[Lindenbaum--Tarski algebra] The formula algebra modulo provable
interderivability; see Definition~\ref{def:p10-lindenbaum-algebra}.
\item[model checking] Determining whether a fixed finite structure/state
satisfies a formula; distinct from satisfiability over all structures.
\item[normal form] A formula satisfying a declared syntactic restriction and
connected to the input by equivalence or equisatisfiability as explicitly stated.
\item[proof sketch] An incomplete proof presentation whose omitted steps remain
owed and which is never cited as a complete proof.
\item[realizability] An interpretation assigning computational witnesses to
formulas under a specified applicative structure.
\item[relative consistency] A metatheoretic implication
\(\operatorname{Con}(T)\to\operatorname{Con}(U)\), not an absolute consistency proof.
\item[saturation] Realization of specified types over parameter sets below a
declared cardinal bound.
\item[semantic consequence] Truth preservation over a declared class of
interpretations; see Definition~\ref{def:p11-model-consequence}.
\item[soundness] The assertion that derivability in a named calculus implies
semantic consequence under a named semantics.
\item[stable model] A fixed point of the Gelfond--Lifschitz reduct construction;
see Definition~\ref{def:p09-stable-model}.
\item[Stone space] The compact zero-dimensional Hausdorff dual of a Boolean
algebra; see Definition~\ref{def:p10-stone-space}.
\item[substitution] Capture-avoiding replacement with the exact term/formula
and free-variable conditions stated.
\item[theory] A set of sentences, often deductively closed only when explicitly
said so.
\item[type] A consistent set of formulas in specified free variables over a
declared parameter set.
\item[ultrafilter] A maximal proper filter; see Definition~\ref{def:p10-boolean-algebra}.
\item[validity] Truth in every interpretation in a declared class.
\item[variety] A class of similar algebras axiomatized by equations;
equivalently an HSP-closed class when arbitrary products include the empty
product, by Theorem~\ref{thm:p10-hsp}.
\end{description}


Within each domain, entries are alphabetized by their leading term. A grouped
entry supplies the linked distinctions rather than scattering near-synonyms.
An exact definition's domain conditions---for example positive conditioning
mass, nonempty action menus or a specified extension family---remain part of
the meaning. The glossary is a cross-domain reference, not a replacement for
the separate source-presentation entries.

\subsection{Logic, syntax, semantics and proof}
\label{sec:p14-glossary-logic}

\begin{description}[style=nextline,leftmargin=0pt]
\item[Alpha-equivalence.]
Consistent renaming of bound variables under the specified grammar and
freshness conditions. It preserves binding, not literal spelling.
\cref{def:p02-fol-alpha-equivalence,def:p07-alpha-substitution}.
\item[Atom; literal.]
An atomic formula belongs to the active grammar's atomic category. In the
categorical fragment, atoms are unary predicate symbols and literals are
those atoms or their formal complements; in first-order syntax a relation
applied to terms, or an equality of terms, is an atomic formula.
\cref{def:p02-syllogistic-language,sec:fol-syntax-semantics-02}.
\item[Argument; countermodel.]
An argument is its premises and conclusion. For sentence arguments a
countermodel satisfies every premise but not the conclusion; with free
variables, the refuting semantic data include an assignment.
\cref{sec:levels-05}.
\item[Calculus; derivation; theorem.]
A calculus inductively specifies initial derivations and rules. Derivability
asserts existence of a licensed derivation with the allowed open assumptions;
a theorem has none. \cref{sec:levels-01,sec:levels-03}.
\item[Compactness.]
The first-order consequence theorem reduces consequence to a finite subset
of premises; the equivalent satisfiability form tests finite subsets of a
theory. Its ambient assumptions and signature scope are stated in
\cref{thm:fol-compactness}. Do not infer an effective decision procedure
from this finite-witness property.
\item[Completeness; model completeness.]
Calculus completeness connects semantic consequence to derivability.
Deductive completeness of a theory decides every sentence by derivation of
it or its negation; its definition does not by itself assert consistency.
Model completeness concerns elementary embeddings between models.
\cref{sec:levels-03,def:theory-completeness}.
\item[Cut; sequent; context.]
A sequent has the active sheet's antecedent and succedent shape, with
contexts declared as sets, multisets or sequences. Cut combines two
derivations through the displayed intermediate formula. Eliminability
means that the same endsequent has a derivation without that rule, not that
the rule is unsound. The exact single-succedent result is
\cref{thm:p07-lj-exact-cut-elimination}; compare \cref{sheet:lj,sheet:lk}.
\item[Elementary equivalence; elementary substructure; isomorphism.]
Elementary equivalence is agreement on sentences. Elementary substructure
preserves formulas with parameters along the inclusion. Isomorphism is
a structure-preserving bijection; it implies agreement on satisfaction
as proved in \cref{lem:iso-invariance}. See \cref{tab:p03-model-symbols}.
\item[Expansion; reduct; substructure.]
Expansion adds interpretations in an extended signature on the same domain;
reduct forgets them. A substructure restricts the domain with closure under
functions and the induced relations. These are not the same operation.
\cref{sec:fol-semantics-01,sec:maps-elementarity-02}.
\item[Formula; sentence; free occurrence.]
A formula belongs to the declared inductive grammar. A sentence has no
free variable occurrences. Free and bound status belongs to occurrences,
not to one letter everywhere in an expression.
\cref{sec:fol-syntax-semantics-02,sec:fol-syntax-semantics-03}.
\item[Language; signature; arity.]
A formal language consists of expressions admitted by its alphabet and
grammar. A first-order signature supplies nonlogical symbols with fixed
arities; arity fixes how many arguments a symbol takes.
\cref{sec:levels-01,sec:fol-syntax-semantics-01}.
\item[Logical equivalence; equisatisfiability.]
Truth-value equivalence uses the same permitted valuation for both formulas;
equisatisfiability merely equates existence of a satisfying valuation or
structure. Equisatisfiable formulas need not be interchangeable within a
larger formula. \cref{thm:p02-replacement,
def:p02-substitution-replacement-limits,thm:skolem-equisatisfiable}.
\item[Model; satisfaction; validity.]
A model satisfies the specified sentences. Satisfaction uses the complete
semantic data, while validity quantifies over all permitted data. Truth in
one structure is not validity over a class.
\cref{sec:levels-04,sec:levels-03,sec:fol-semantics-02}.
\item[Normal form; normalization; confluence.]
A beta-normal term has no beta redex. Weak normalization asserts existence
of a path to a normal form; strong normalization forbids infinite reduction.
Confluence joins divergent finite reduction paths; it does not assert that
either path terminates. \cref{def:p07-beta-conversion,
def:p07-confluence-not-normalization,thm:p07-beta-normal-form-unique}.
Formula normal-form transformations separately state equivalence or
equisatisfiability; see \cref{thm:prenex-normal-form,
thm:skolem-equisatisfiable}.
\item[Object language; metalanguage.]
The language studied and the language used to define and analyze it.
A formula inscription, a statement about the inscription and a statement
about its code have different roles. \cref{sec:levels-02}.
\item[Paraconsistency; explosion; dialetheism.]
Paraconsistency is failure of explosion for a consequence relation: a
contradiction does not entail every conclusion. Dialetheism is the further
thesis that some contradiction is true. The former alone does not assert
the latter. \cref{def:p08-paraconsistency}.
\item[Semantic consequence; soundness.]
Consequence preserves satisfaction over the declared semantic class.
Soundness sends derivability to consequence, with language, calculus and
semantics fixed. \cref{sec:levels-03,thm:fol-soundness}.
\item[Skolemization.]
Replace existential witnesses by fresh function symbols with the stated
dependencies and retain the universal prefix. Its equisatisfiability theorem
compares expanded signatures; it does not assert literal equivalence in
the original language. \cref{thm:skolem-equisatisfiable}.
\item[Substitution; selected replacement.]
Use the grammar-specific operations of
\cref{def:p02-uniform-substitution,
def:p02-fol-capture-avoiding-substitution,def:p02-selected-replacement}.
Uniformity, simultaneous insertion, freedom from capture and selected
occurrence positions are separate conditions.
\item[Theory; type.]
A theory is a set of sentences; deductive closure is an additional
convention when imposed. A model-theoretic type is a set of formulas in
specified variables consistent with the declared parameter structure or
theory. A complete type chooses each formula or its negation, and need
not be realized in a chosen model. \cref{def:fol-theory-satisfiability,
def:type-space}.
\end{description}

\subsection{Inference, explanation, computation and argumentation}
\label{sec:p14-glossary-inference}

\begin{description}[style=nextline,leftmargin=0pt]
\item[Abduction; abducibles; manifestations.]
In the finite logical entry, abducibles are permitted hypothesis atoms and
manifestations are observation atoms. An explanation is a permitted set
jointly satisfiable with the background and sufficient to entail the
observations. Neither sufficiency nor minimality proves its truth.
\cref{def:p16-abduction-input,def:p16-abduction-models,
thm:p16-abduction-truth-boundary}.
\item[Admissibility; defense.]
An argumentation set defends an argument when it attacks each attacker of
that argument. It is admissible when conflict-free and defending all its
own members. Defense does not require membership.
\cref{def:p18-defense-operator}.
\item[Attack; conflict-free set.]
An attack is an edge in the supplied directed argumentation relation. A
conflict-free set contains no internal attack, including a self-attack.
The graph does not certify premise truth or inference validity.
\cref{def:p18-finite-framework}.
\item[Computably enumerable; decidable; recognition.]
A c.e. numerical set is the domain of a partial computable function.
A decidable set has a total computable characteristic function. Recognition
accepts exactly the members but may diverge on nonmembers; decision halts
on every permitted input. \cref{sec:turing-machines-03,
thm:machine-decidability,thm:ce-equivalent-forms}.
\item[Credulous; skeptical acceptance.]
Credulous acceptance is membership in some selected extension; skeptical
acceptance here requires a nonempty extension family and membership in
every one of its members. No extensions is distinct from the single empty
extension. \cref{def:p18-acceptance-queries}.
\item[Deduction; inductive support.]
Deduction preserves truth under its fixed premises and semantics. Statistical
or empirical support uses its specified data, comparison and model, and
need not entail the hypothesis. Their outputs and failure witnesses are
compared in \cref{sec:p14-evidence-support,
def:p19-evidence-record}; a valid argument can still have false premises.
\item[Defeasibility; nonmonotonicity.]
Added information can withdraw a conclusion under the selected consequence
relation. Smooth preferential models select minimal premise states;
ranking imposes a further layered structure.
\cref{def:p08-nonmonotonic,def:p08-preferential-model}.
\item[Extension: complete, grounded, preferred, stable.]
A complete extension is admissible and includes every argument it defends;
the grounded extension is the least fixed point of defense; a preferred
extension is inclusion-maximal admissible; a stable extension is
conflict-free and attacks every outside argument. Maximal does not mean
maximum cardinality. \cref{def:p18-extension-semantics}.
\item[Heuristic; admissible; consistent.]
A search heuristic estimates remaining path cost; admissibility makes it
a lower bound, and consistency imposes the edgewise cost inequality.
The graph/interface conditions still determine the algorithmic guarantee.
\cref{def:p12-hsearch-instance,def:p12-hsearch-heuristics,
thm:p12-hsearch-consistency-bound}.
\item[Minimal; minimum.]
A subset-minimal explanation has no proper explanatory subset. A
minimum-cardinality explanation has as few hypotheses as any explanation.
Different incomparable minimal explanations can survive.
\cref{def:p16-abduction-minimality,thm:p16-abduction-minimality}.
\item[Model checking; satisfiability.]
Model checking tests a supplied structure/state against a formula;
satisfiability asks whether some permitted semantic data satisfy it.
Algorithmic complexity depends on which objects are inputs and their
encoding. \cref{sec:p14-decision-table,
def:fol-theory-satisfiability}.
\item[Monitoring; control; feeling of rightness.]
Monitoring judges the current reasoning activity; control changes its
strategy, resources or stopping decision. Feeling of rightness accompanies
an initial answer and is not automatically a probability report.
\cref{def:p19-r86-monitoring-control,def:p19-r86-judgment-types}.
\item[Normative benchmark; empirical evidence.]
A benchmark states the standard and its assumptions. Evidence records
responses from a population/task/design and constrains a claim about that
record; matching an answer key does not identify a unique cognitive process.
\cref{def:p19-benchmark,def:p19-empirical-question,
def:p19-evidence-record,sec:p14-evidence-support}.
\item[Resolution; refutation.]
The first-order resolution method operates on the stated clause language
with its substitution and lifting conditions. Its completeness theorem
concerns refutation of unsatisfiable clause inputs, not a universal
termination guarantee for every search. \cref{thm:resolution-lifting,
thm:fol-resolution}.
\end{description}

\subsection{Probability, statistics and uncertainty}
\label{sec:p14-glossary-probability}

\begin{description}[style=nextline,leftmargin=0pt]
\item[Almost surely; null event; support.]
Probability one, probability zero and the positive-mass outcomes,
respectively, in the finite model. A nonempty event may be null; an
almost-sure property need not hold at every outcome.
\cref{def:p15-finite-probability-space,def:p15-null-conditioning}.
\item[Bayes factor; posterior model odds.]
A Bayes factor is the ratio of two marginal data densities, integrating
each model's proper parameter prior on the stated domain. Posterior odds
multiply it by model prior odds; a factor is not itself a posterior
probability. \cref{def:p15-model-comparison-experiment,
thm:p15-model-posterior-odds}.
\item[Belief function; plausibility; possibility; necessity.]
Finite belief and plausibility sum subset masses by containment or overlap;
possibility takes a maximum of singleton scores, and necessity is its
complement dual. Their exact union laws differ from additive probability.
\cref{def:p14-evidence-uncertainty-data,
thm:p14-evidence-finite-measures,thm:p14-evidence-measure-separation}.
\item[Calibration; accuracy; resolution.]
Calibration equates each positive-mass reported probability with its
conditional outcome frequency under the specified joint law. Accuracy
is mean correctness for a stated key; resolution varies outcome rates
between judgment classes. None is another name for either of the others.
\cref{def:p19-calibration-law,thm:p19-calibration-decomposition}.
\item[Common effect; random effects; heterogeneity.]
A common-effect synthesis asserts one study estimand. Random-effects
synthesis posits a distribution of underlying study effects; its mean and
variance differ from within-study sampling uncertainty and from a
deterministically weighted finite target. Systematic bias is not made
mean zero by the name random effect.
\cref{def:p15-synthesis-record,def:p15-synthesis-random-model}.
\item[Conditional probability; conditional expectation.]
Event conditioning reweights by a positive event probability. Finite
conditional expectation averages on the defined conditioning strata.
Neither operation by itself changes structural equations.
\cref{def:p15-finite-conditioning,
def:p15-finite-conditional-expectation}.
\item[Confidence set; credible set.]
Sampling coverage includes the fixed target with the claimed probability
for every parameter in scope. Posterior content is conditional probability
of membership after the data. The separate prior-average coverage identity
does not turn posterior content into uniform fixed-parameter coverage.
\cref{def:p15-confidence-procedure,def:p15-bayes-credible-set,
thm:p15-bayes-coverage}.
\item[Covariance; correlation; uncorrelated.]
For variables with finite second moments, covariance averages the product
of their centered values. Correlation divides covariance by their two
standard deviations and is undefined by that formula if either variance
is zero. Uncorrelated means zero covariance, not the independence
factorization. \cref{def:p15-covariance-correlation}.
\item[Cumulative distribution function; quantile.]
A cumulative distribution function gives the probability of a real
random variable being at most a threshold. The quantile map uses the
specified generalized inverse; it does not require an ordinary inverse
or a strictly increasing continuous cumulative function.
\cref{def:p15-cumulative-distribution,def:p15-quantile-map}.
\item[Density; mass function; law.]
A finite mass function weights individual outcomes and totals one. A
random variable's law groups those outcomes by value; a density represents
a law relative to a specified reference measure. A density value at one
point need not be a point probability.
\cref{def:p15-finite-probability-space,def:p15-finite-random-variable,
def:p15-bayes-experiment}.
\item[Estimate; estimand; analysis sample.]
The estimand fixes the quantity, population, task and comparison. The
estimate is computed by the stated procedure on the observations; the
analysis sample reflects the recorded exclusion and missing-data rules.
\cref{def:p19-population-and-estimand,def:p19-evidence-record}.
\item[Event; sigma-algebra; probability space.]
An event is a member of the declared measurable collection. A sigma-algebra
is closed under complements and countable unions; a probability measure
is countably additive there and gives the whole space mass one. A finite
space with all subsets measurable is a special case, not the definition
for every infinite space. \cref{def:p15-sigma-algebra,def:p15-measure,
def:p15-finite-probability-space}.
\item[Expectation.]
In the finite setting, the probability-weighted sum of numerical values;
linearity concerns the expectation operation and does not require
independence. An expectation does not describe every realized outcome.
\cref{def:p15-finite-expectation}.
\item[Independence: pairwise, mutual, conditional.]
Product factorization for pairs, for every subfamily, or under the
specified conditional law. Pairwise independence alone does not imply
mutual independence, and selection can change conditional relationships.
\cref{def:p15-finite-independence,def:p15-pairwise-not-mutual}.
\item[Likelihood; likelihood ratio.]
At fixed observed data, likelihood is a sampling density regarded as a
function of its parameter. It need not normalize over parameters. A ratio
must state numerator, denominator and zero-denominator treatment.
\cref{def:p15-bayes-experiment,def:p15-likelihood-test}.
\item[Lower probability; lower prevision; coherence.]
An event lower envelope ranges over a retained nonempty family of laws.
A lower prevision specifies supremum acceptable buying prices for
bounded gambles; coherence includes its stated portfolio conditions.
Event envelopes and gamble assessments have different domains.
\cref{def:p14-evidence-uncertainty-data,def:p15-ip-domain,
def:p15-ip-prices,def:p15-ip-coherence}.
\item[Measurement uncertainty; expanded uncertainty.]
A standard uncertainty describes the selected uncertainty model's scale;
expanded uncertainty multiplies the combined standard uncertainty by
a stated coverage factor. The multiplier alone supplies no universal
coverage probability. \cref{def:p15-measurement-expanded,
thm:p15-measurement-coverage-models}.
\item[Odds; prior; posterior.]
Odds compare the probability of an event with its complement. A prior is
part of the stated joint model; a posterior conditions that model on the
observed evidence on its normalization domain.
\cref{def:p15-finite-hypothesis-partition,thm:p15-finite-bayes,
def:p15-bayes-experiment}.
\item[\(p\)-value; level; size; power.]
A valid \(p\)-value is superuniform under each null law. Size is the
supremum null rejection probability; a level bounds size; power is
rejection probability at a specified alternative. None is the
probability that the null is true.
\cref{def:p15-valid-p-value,def:p15-statistical-experiment}.
\item[Prediction; prediction interval.]
A predictive law concerns a specified future observation or effect,
not merely an unknown mean. Random-effects prediction includes the
new-effect variability; replacing known variances by estimates needs
its separately stated approximation and uncertainty treatment.
\cref{def:p15-bayes-future-kernel,thm:p15-bayes-predictive-kernel,
thm:p15-synthesis-random-known,thm:p15-synthesis-estimated-weights}.
\item[Random variable; joint law; marginal law.]
A finite random variable is a function on the outcome space. Its joint
law with another variable records pairs; the marginal law sums the joint
masses over the other coordinate. The outcome space and the value space
are not automatically the same set.
\cref{def:p15-finite-random-variable,thm:p15-finite-law-projection}.
In the general setting a random element is a measurable map into the
specified value space: inverse images of measurable value sets must be
events. Its law is the pushforward probability measure.
\cref{def:p15-measurable-map,def:p15-pushforward-law}.
\item[Type-I error; type-II error.]
Rejecting at a null law and failing to reject at an alternative law,
respectively. The probabilities refer to the specified test and law;
nonrejection does not prove the null.
\cref{def:p15-statistical-experiment}.
\end{description}

\subsection{Causation and counterfactual reasoning}
\label{sec:p14-glossary-causation}

\begin{description}[style=nextline,leftmargin=0pt]
\item[Abduction, action, prediction.]
For a specified finite causal model and positive factual evidence,
abduction reweights the background outcomes, action replaces queried
equations, and prediction evaluates the modified worlds under the same
posterior background. This use of abduction is not the finite logical
explanation search. \cref{def:p16-counterfactual-evidence,
thm:p16-counterfactual-calculation}.
\item[Average treatment effect; effect among the treated.]
ATE averages \(Y(1)-Y(0)\) over the specified population; ATT averages it
conditional on received treatment one. Their target weights and
identification assumptions differ.
\cref{def:p16-binary-potential-outcomes,
thm:p16-po-ate-identification,thm:p16-po-att-identification}.
\item[Back-door criterion; adjustment.]
The specified covariates contain no treatment descendants and block
every back-door path in the complete DAG. The overlap condition is an
additional requirement for the conditional ratios used by the finite
adjustment calculation. \cref{def:p16-back-door-data}.
\item[Collider; d-separation.]
Collider status concerns the two incoming edges used on a path. A path
is active when its internal noncolliders are unconditioned and its
colliders have themselves or a descendant conditioned on. D-separation
is absence of such active paths between the disjoint endpoint sets.
\cref{def:p16-active-path}. It is a graph relation, distinct from
distributional conditional independence in
\cref{def:p16-graphical-conditional-independence}.
\item[Counterfactual world; twin construction.]
Worlds use specified equation replacements with one shared background
outcome; the twin construction duplicates endogenous coordinates while
retaining their common background. Independently sampling a new background
for each world changes the joint query.
\cref{def:p16-counterfactual-worlds,def:p16-twin-data,
thm:p16-finite-twin-construction}.
\item[Endogenous; exogenous.]
Endogenous variables receive values from the model's indexed structural
equations. Exogenous disturbances have a supplied joint law but no
endogenous assignment equations in that model. Exogenous does not mean
that the disturbances are independent.
\cref{def:p16-finite-scm}.
\item[Exchangeability; consistency; positivity.]
In the binary-treatment entry, conditional exchangeability is independence
of each potential outcome and received treatment given the covariates;
consistency identifies the observed outcome with the received-treatment
potential outcome; positivity requires the relevant treatment probability
in the covariate strata used. Identification uses their stated combination,
not any one of them alone.
\cref{def:p16-binary-potential-outcomes,thm:p16-po-ate-identification}.
\item[Identification; estimation; nonidentification.]
Identification makes the target equal across every allowed model with
the same observational law. A pair sharing that law but differing in
the target witnesses nonidentification. Estimation from sampled data
remains a further problem even when identification holds.
\cref{def:p16-causal-identification,thm:p16-causal-nonidentification}.
\item[Intervention; observation.]
Intervention replaces the specified equations while preserving the other
mechanisms and exogenous law as stipulated. Observation conditions on an
event under unchanged equations. The do-expression is an operation, not
an event requiring positive observational probability.
\cref{def:p16-causal-intervention,def:p15-finite-conditioning}.
\item[Potential outcome; probability of benefit.]
Potential outcomes give responses under the stated treatment interventions
on a common background. The probability of benefit depends on their joint
coupling; equal marginal intervention laws do not generally identify it.
\cref{def:p16-binary-potential-outcomes,
thm:p16-po-coupling-nonidentification}.
\item[Structural causal model.]
The finite acyclic entry supplies a DAG, value sets, a background law,
disturbances and total indexed structural functions. Its permitted parent
edge need not imply a nonzero effect or observed association.
\cref{def:p16-causal-dag,def:p16-finite-scm}.
\end{description}

\subsection{Decision, preference and strategic interaction}
\label{sec:p14-glossary-decision}

\begin{description}[style=nextline,leftmargin=0pt]
\item[Act; lottery; mixed strategy.]
A finite lottery assigns nonnegative prize probabilities totaling one.
A lottery act maps a state to a lottery; a mixed strategy is a lottery
on a player's pure actions. The spaces and consequence interpretation
must be supplied. \cref{def:p17-finite-lottery,
def:p17-ambiguity-acts,def:p17-finite-strategic-game}.
\item[Bayes action; posterior risk; loss.]
Posterior risk averages a supplied loss under the posterior. A Bayes
action minimizes that risk among the feasible actions; the posterior
does not select the loss function or menu.
\cref{def:p15-bayes-loss,thm:p15-bayes-decision-rule}.
\item[Best response; Nash equilibrium.]
A best response maximizes one player's expected payoff with the opponents'
strategies fixed. A Nash profile permits no profitable unilateral mixed
deviation; this condition is not a claim of collective optimality or an
algorithmic convergence guarantee.
\cref{def:p17-nash-equilibrium,thm:p17-best-response-geometry}.
\item[Expected utility; consequence utility.]
Consequence utility assigns prize values; expected utility weights them
by lottery probabilities. The representation needs its stated preference
axioms, and positive affine scale changes preserve the lottery ordering.
\cref{def:p17-expected-utility,thm:p17-finite-vnm,
thm:p17-vnm-affine-uniqueness}.
\item[Independence axiom.]
Mixing each of two lotteries with the same third lottery at the same
weight strictly between zero and one preserves and reflects their preference
comparison. It is not event independence and deliberately excludes the
zero-weight endpoint. \cref{def:p17-lottery-independence}.
\item[Indifference; strict preference; weak preference.]
Weak preference is the declared complete transitive relation. Strict
preference is one affirmative weak comparison with the reverse absent;
indifference has both. Lack of comparison is not indifference.
\cref{def:p17-weak-preference,def:p17-preference-quotient}.
\item[Maximal choice; maximizer.]
The finite weak-preference choice correspondence returns alternatives
weakly preferred to every available alternative. In more general partial
orders, inclusion-maximality or nondomination does not mean greatest
numerical value or maximum cardinality. Use the exact relation and
menu in \cref{def:p17-maximal-choice,
def:p18-extension-semantics}.
\item[Maxmin expected utility; robust dominance.]
Maxmin ranks an act by its worst expected utility over the retained
priors. Robust dominance requires the comparison under every retained
prior. These can give different orderings, and the prior list is an
input, not uniquely inferred by naming ambiguity.
\cref{def:p17-finite-maxmin,def:p17-robust-dominance}.
\item[Ordinal representation; cardinal lottery utility.]
An ordinal representation preserves the ordering of sure alternatives;
increasing transformations of attained values preserve it. Lottery
expected-utility representations restrict permissible utility scale
changes to positive affine transformations. An ordinal prize ranking
alone does not fix uncertain comparisons.
\cref{def:p17-ordinal-representation,
def:p17-increasing-scale-change,thm:p17-vnm-affine-uniqueness}.
\item[Policy; Markov policy; stationary policy.]
A sequential policy chooses a distribution on feasible actions after
each available history. A Markov policy depends on the current time and
state; a deterministic stationary policy in a time-homogeneous model
uses one feasible action map on states. These are restrictions on the
information and timing interface, not a general claim that the observed
environment is Markovian. \cref{def:p17-sequential-model}.
\item[Pure strategy; strategic dominance.]
A pure strategy is an available action. Strict strategic domination
compares one player's payoff against every opponent profile under the
chapter's condition; it is not the same as a best response against one
fixed profile. \cref{def:p17-finite-strategic-game,
def:p17-strategic-dominance}.
\item[Risk aversion; risk neutrality.]
For the stated interval and finite-support lottery preferences, weak risk
aversion prefers the sure mean to every lottery with that mean; neutrality
is indifference in all those comparisons. Strict preference for every
nonconstant lottery's sure mean is a stronger property. These are
comparisons under a supplied consequence utility, not classifications
deduced merely from choosing a lower-variance option once.
\cref{def:p17-concavity-risk}.
\item[Stochastic dominance: first order.]
The finite-support entry compares expected utilities for every continuous
strictly increasing consequence utility on the stated interval. It is
distinct from strategic domination and from comparison under one fixed
utility function. \cref{def:p17-first-order-dominance,
def:p17-strategic-dominance}.
\item[Value of information; perfect information.]
The finite signal model evaluates the best policy conditional on the
available signal. No-information value uses one unconditional action;
perfect-information value observes the state itself. Their difference
depends on the prior, payoffs, feasible actions and signal experiment,
not the number of facts supplied alone.
\cref{def:p17-information-policy,thm:p17-information-value-bounds}.
\end{description}

\section{Cross-references between standard names and historical terminology}
\label{sec:p14-glossary-names}

A historical name is a retrieval key, not evidence that every later version
has the original grammar, premises or proof. Use the source-presentation
entry for an original formulation and the linked handbook theorem for the
stated modern version. The biography and chronology in
\cref{sec:p13-biographies,sec:p13-chronology} supply historical context;
the named-theorem index supplies statement and proof destinations.
The notation histories in \cref{sec:p01-notation-variants,
sec:p13-notation-history} record coexistence and variation, not a sequence
in which one notation made all its predecessors obsolete
\cite{WhiteheadRussell1910,Lukasiewicz1929,HilbertAckermann1928,
Gentzen1934a,Gentzen1935b}.

\begin{longtable}{@{}p{.24\textwidth}p{.30\textwidth}p{.36\textwidth}@{}}
\caption{Aliases, source terms and qualifications for historical lookup.}
\label{tab:p14-glossary-names}\\
\toprule
Name encountered & Current lookup or source meaning & Qualification and destination \\
\midrule
\endfirsthead
\caption[]{Aliases, source terms and qualifications for historical lookup (continued)}\\
\toprule
Name encountered & Current lookup or source meaning & Qualification and destination \\
\midrule
\endhead
\midrule\multicolumn{3}{r}{\emph{continued on next page}}\\
\endfoot
\endlastfoot
Alternative denial; Sheffer function; NAND & Boolean output 1110 in the
declared 00,01,10,11 input order. & The historical stroke is not a global
active operator; read the defined function and its source-specific glyph.
\cref{tab:p01-sixteen,sec:p01-notation-variants}. \\
Joint denial; Peirce arrow; NOR & Boolean output 1000 in that input order. &
Neither inclusive nor exclusive disjunction. \cref{tab:p01-sixteen}. \\
Material implication; conditional & The truth function false only when the
antecedent is true and consequent false. & Historical \(\supset\),
Polish \(C\) and modern \(\to\) compare local spellings, not all causal,
counterfactual or strict conditionals. \cref{tab:p01-sixteen,
tab:p01-notation-variants}. \\
Equivalence; biconditional & Boolean biconditional in the connective
position; metalinguistic formula equivalence in another position. &
Historical \(\equiv\) can be a source connective, while structure
operands use elementary equivalence. \cref{tab:p01-notation-variants,
thm:p02-replacement,tab:p03-model-symbols}. \\
Polish notation & Prefix notation in the identified {\L}ukasiewicz
presentation, with the connective before its operands. & Letter \(A\)
can name disjunction there rather than a formula metavariable; identify
grammar and arity before parsing. \cref{tab:p01-notation-variants}. \\
Universal categorical; existential categorical & All/no/some categorical
class claims, represented by the two-literal forms of the modern fragment. &
The modern empty-class policy is not automatically an attribution to
Aristotle. \cref{def:p02-syllogistic-language,
def:p02-syllogistic-semantics,sec:p02-syllogistic-inference}. \\
Hauptsatz & Gentzen's cut-elimination result and its specified
applications. & Exact modern sheets keep their own sequent shape and
negation conventions; the single-succedent LJ result is
\cref{thm:p07-lj-exact-cut-elimination}. Original and later host proofs
must not be identified solely by this name. \cite{Gentzen1934a,Gentzen1935b}. \\
Entscheidungsproblem; decision problem & A specified effective yes/no
problem, classically first-order validity in the historical context. &
Undecidability is about the named input problem, not an inability to reason
about every individual case. \cref{thm:fol-validity-undecidable,
sec:p14-decision-table} \cite{Church1936,Turing1937}. \\
Recursive set; decidable set & The set has a total computable characteristic
function under the fixed coding. & Recursive numerical functions and
recursive definitions have their own type and totality requirements.
\cref{sec:turing-machines-03,thm:machine-decidability}
\cite{Shoenfield1967}. \\
Recursively enumerable; r.e.; computably enumerable; c.e. & The domain of
a partial computable function; equivalently recognizable under the stated
machine presentation. & Enumeration does not assert decision of
nonmembership. Preserve the empty-set exception in the total-range form.
\cref{sec:turing-machines-03,thm:ce-equivalent-forms}
\cite{Shoenfield1967}. \\
Semidecision; recognition & Acceptance of exactly the members with possible
divergence outside the set. & Distinct from a decider's halting on every
permitted input. \cref{sec:turing-machines-03}. \\
Skolem function; Skolemization & A fresh function supplies existential
witnesses with the stated universal dependencies. & Equisatisfiability
uses an expanded signature and retains the stated Choice boundary.
\cref{thm:skolem-equisatisfiable}. \\
L\"owenheim--Skolem; Skolem paradox & The cardinality theorem and the
internal/external countability distinction discussed in the historical
account. & Not a contradiction showing that a model proves its own
uncountable set countable. The biography gives the contextual distinction;
use the actual cardinality theorem's hypotheses.
\cref{sec:p13-biographies,sec:fol-cardinality-normalforms-01}. \\
Manifestations; abducibles & Eiter--Gottlob observation atoms and permitted
hypothesis atoms. & The linked finite explanation criterion is not a
posterior or a causal-effect definition.
\cref{def:p16-abduction-input,def:p16-abduction-models}. \\
Acceptable with respect to a set & Dung's defense condition. & It does not
assert membership, factual truth or skeptical acceptance.
\cref{def:p18-defense-operator,def:p18-acceptance-queries}. \\
Distribution function, in the finite CHANCE source & The source's
atom-weight function, called a mass function in the finite handbook entry. &
Not the cumulative distribution function or a density.
\cref{def:p15-finite-probability-space}
\cite{GrinsteadSnellChance2006}. \\
Confidence coefficient, in Neyman 1937 & The procedure's claimed
sampling coverage, written \(1-\alpha\) in the current chapter. &
Neyman's coefficient notation must not be read as this chapter's
significance level. \cref{def:p15-confidence-procedure}
\cite{Neyman1937Confidence}. \\
Bayes' formula; Bayes factor & Conditional inversion in a joint model;
integrated model-likelihood ratio, respectively. & Shared attribution
does not make the two numerical objects equal.
\cref{thm:p15-finite-bayes,thm:p15-model-posterior-odds}. \\
Fixed-effect meta-analysis; common-effect target & In the common-target
entry, every study estimates the same underlying effect. & Always state
whether a source instead means a finite weighted target; this differs
from a random-effects population mean.
\cref{def:p15-synthesis-record}
\cite{DeeksEtAl2024Cochrane10}. \\
Bernoulli utility; von Neumann--Morgenstern utility & Utility of consequences
in the finite expected-utility representation. & Distinct from its
expected-utility function on lotteries and from an arbitrary ordinal
representation of sure alternatives.
\cref{def:p17-expected-utility,thm:p17-vnm-affine-uniqueness}. \\
Countering; best-response correspondence & The correspondence used by
Nash's original equilibrium existence route. & Its fixed-point condition
concerns unilateral deviation, not a collective payoff maximum.
\cref{def:p17-nash-equilibrium,thm:p17-best-response-geometry}
\cite{Nash1950EquilibriumPoints,Nash1951NonCooperative}. \\
Meta-reasoning; metacognitive monitoring and control & Judgments of the
reasoning activity and actions changing that activity. & Psychological
object/meta levels are not automatically formal object-language and
metalanguage levels. \cref{def:p19-r86-monitoring-control}
\cite{ThompsonElqayamAckerman2021MetaReasoning}. \\
\bottomrule
\end{longtable}

When searching by a person's name, also search the distinctive result and its
input class. A theorem's attributed name, a calculus carrying that name and
a later algorithm using the theorem can have different statements. In
particular, a soundness theorem is not a completeness theorem, a completeness
theorem is not a decision algorithm, and a decision algorithm does not settle
whether its input formalization expresses the intended reasoning task.
The task-based routes in \cref{sec:p14-reasoning-task-finder} and the
boundary cases in \cref{sec:p14-countermodels-and-boundary-cases} help select
the required destination.
